\documentclass[10pt,a4paper]{amsart}
\usepackage{amsmath,amssymb,amsthm,graphicx}
% Retained preamble for the Stacks Project chapter snapshot in this directory.
% The chapter's own preamble.tex is a copy of the shared Stacks Project preamble; the block
% below is sliced MECHANICALLY, line for line, from it. Nothing was added, removed or
% redefined; the wrapper class, the xr-hyper/hyperref wiring and the externaldocument list
% are omitted because the assembler supplies the document class.
%
% lines taken from preamble.tex: 3, 5, 9, 11, 12, 156, 157, 158, ... (51 lines total)
%
% package coverage restored from the chapter's own preamble (2026-10-13): xy, xy-2cell

\usepackage[all]{xy}
\xyoption{2cell}
\UseAllTwocells

\newenvironment{reference}{\comment}{\endcomment}
\newenvironment{slogan}{\comment}{\endcomment}
\newenvironment{history}{\comment}{\endcomment}
\theoremstyle{plain}
\newtheorem{theorem}[subsection]{Theorem}
\newtheorem{proposition}[subsection]{Proposition}
\newtheorem{lemma}[subsection]{Lemma}
\theoremstyle{definition}
\newtheorem{definition}[subsection]{Definition}
\newtheorem{example}[subsection]{Example}
\newtheorem{exercise}[subsection]{Exercise}
\newtheorem{situation}[subsection]{Situation}
\theoremstyle{remark}
\newtheorem{remark}[subsection]{Remark}
\newtheorem{remarks}[subsection]{Remarks}
\numberwithin{equation}{subsection}
\def\lim{\mathop{\mathrm{lim}}\nolimits}
\def\colim{\mathop{\mathrm{colim}}\nolimits}
\def\Spec{\mathop{\mathrm{Spec}}}
\def\Hom{\mathop{\mathrm{Hom}}\nolimits}
\def\Ext{\mathop{\mathrm{Ext}}\nolimits}
\def\SheafHom{\mathop{\mathcal{H}\!\mathit{om}}\nolimits}
\def\SheafExt{\mathop{\mathcal{E}\!\mathit{xt}}\nolimits}
\def\Sch{\mathit{Sch}}
\def\Mor{\mathop{\mathrm{Mor}}\nolimits}
\def\Ob{\mathop{\mathrm{Ob}}\nolimits}
\def\Sh{\mathop{\mathit{Sh}}\nolimits}
\def\NL{\mathop{N\!L}\nolimits}
\def\CH{\mathop{\mathrm{CH}}\nolimits}
\def\proetale{{pro\text{-}\acute{e}tale}}
\def\etale{{\acute{e}tale}}
\def\QCoh{\mathit{QCoh}}
\def\Ker{\mathop{\mathrm{Ker}}}
\def\Im{\mathop{\mathrm{Im}}}
\def\Coker{\mathop{\mathrm{Coker}}}
\def\Coim{\mathop{\mathrm{Coim}}}
\DeclareMathSymbol{\boxtimes}{\mathbin}{AMSa}{"02}
\def\QCohstack{\mathcal{QC}\!\mathit{oh}}
\def\Cohstack{\mathcal{C}\!\mathit{oh}}
\def\Spacesstack{\mathcal{S}\!\mathit{paces}}
\def\Quotfunctor{\mathrm{Quot}}
\def\Hilbfunctor{\mathrm{Hilb}}
\def\Curvesstack{\mathcal{C}\!\mathit{urves}}
\def\Polarizedstack{\mathcal{P}\!\mathit{olarized}}
\def\Complexesstack{\mathcal{C}\!\mathit{omplexes}}
\def\Pic{\mathop{\mathrm{Pic}}\nolimits}
\def\Picardstack{\mathcal{P}\!\mathit{ic}}
\def\Picardfunctor{\mathrm{Pic}}
\def\Deformationcategory{\mathcal{D}\!\mathit{ef}}
\usepackage{cleveref}
\crefname{originallemma}{original lemma}{original lemmas}
\Crefname{originallemma}{Original lemma}{Original lemmas}
\makeatletter
\newlabel{algebraic-lemma-characterize-representable-by-algebraic-spaces}{{04Y5}{source}{}{sourcepointer-1}{}}
\newlabel{algebraic-lemma-characterize-representable-by-algebraic-spaces@cref}{{[originallemma][0][]04Y5}{[1][0][]source}{}{}{}}
\makeatother
\crefname{originallemma}{original lemma}{original lemmas}
\Crefname{originallemma}{Original lemma}{Original lemmas}
\makeatletter
\newlabel{algebraic-lemma-equivalent}{{03YQ}{source}{}{sourcepointer-2}{}}
\newlabel{algebraic-lemma-equivalent@cref}{{[originallemma][0][]03YQ}{[1][0][]source}{}{}{}}
\makeatother
\crefname{originallemma}{original lemma}{original lemmas}
\Crefname{originallemma}{Original lemma}{Original lemmas}
\makeatletter
\newlabel{algebraic-lemma-map-space-into-stack}{{04T4}{source}{}{sourcepointer-3}{}}
\newlabel{algebraic-lemma-map-space-into-stack@cref}{{[originallemma][0][]04T4}{[1][0][]source}{}{}{}}
\makeatother
\crefname{originallemma}{original lemma}{original lemmas}
\Crefname{originallemma}{Original lemma}{Original lemmas}
\makeatletter
\newlabel{criteria-lemma-flat-quotient-flat-presentation}{{06FH}{source}{}{sourcepointer-4}{}}
\newlabel{criteria-lemma-flat-quotient-flat-presentation@cref}{{[originallemma][0][]06FH}{[1][0][]source}{}{}{}}
\makeatother
\crefname{originallemma}{original lemma}{original lemmas}
\Crefname{originallemma}{Original lemma}{Original lemmas}
\makeatletter
\newlabel{criteria-theorem-flat-groupoid-gives-algebraic-stack}{{06FI}{source}{}{sourcepointer-5}{}}
\newlabel{criteria-theorem-flat-groupoid-gives-algebraic-stack@cref}{{[originallemma][0][]06FI}{[1][0][]source}{}{}{}}
\makeatother
\crefname{originallemma}{original lemma}{original lemmas}
\Crefname{originallemma}{Original lemma}{Original lemmas}
\makeatletter
\newlabel{lemma-monomorphism}{{lemma-monomorphism}{source}{}{sourcepointer-6}{}}
\newlabel{lemma-monomorphism@cref}{{[originallemma][0][]lemma-monomorphism}{[1][0][]source}{}{}{}}
\makeatother
\crefname{originallemma}{original lemma}{original lemmas}
\Crefname{originallemma}{Original lemma}{Original lemmas}
\makeatletter
\newlabel{lemma-monomorphism-into-point}{{lemma-monomorphism-into-point}{source}{}{sourcepointer-7}{}}
\newlabel{lemma-monomorphism-into-point@cref}{{[originallemma][0][]lemma-monomorphism-into-point}{[1][0][]source}{}{}{}}
\makeatother
\crefname{originallemma}{original lemma}{original lemmas}
\Crefname{originallemma}{Original lemma}{Original lemmas}
\makeatletter
\newlabel{lemma-points-cartesian}{{lemma-points-cartesian}{source}{}{sourcepointer-8}{}}
\newlabel{lemma-points-cartesian@cref}{{[originallemma][0][]lemma-points-cartesian}{[1][0][]source}{}{}{}}
\makeatother
\crefname{originallemma}{original lemma}{original lemmas}
\Crefname{originallemma}{Original lemma}{Original lemmas}
\makeatletter
\newlabel{lemma-unique-point}{{lemma-unique-point}{source}{}{sourcepointer-9}{}}
\newlabel{lemma-unique-point@cref}{{[originallemma][0][]lemma-unique-point}{[1][0][]source}{}{}{}}
\makeatother
\crefname{originallemma}{original lemma}{original lemmas}
\Crefname{originallemma}{Original lemma}{Original lemmas}
\makeatletter
\newlabel{lemma-unique-point-better}{{lemma-unique-point-better}{source}{}{sourcepointer-10}{}}
\newlabel{lemma-unique-point-better@cref}{{[originallemma][0][]lemma-unique-point-better}{[1][0][]source}{}{}{}}
\makeatother
\crefname{originallemma}{original lemma}{original lemmas}
\Crefname{originallemma}{Original lemma}{Original lemmas}
\makeatletter
\newlabel{spaces-morphisms-lemma-noetherian-finite-type-finite-presentation}{{06G4}{source}{}{sourcepointer-11}{}}
\newlabel{spaces-morphisms-lemma-noetherian-finite-type-finite-presentation@cref}{{[originallemma][0][]06G4}{[1][0][]source}{}{}{}}
\makeatother
\title{A unique-point stack contains a unique strictly full reduced locally Noetherian one-point substack}
\author{The Stacks Project Authors}
\date{stacks-project:stacks-properties.tex:local-snapshot-20261012; GNU FDL 1.2 or later}
\begin{document}
\maketitle
\noindent\textit{Editorial note.} Statement span is lines 2790-2802 of sources/stacks-stacks-properties/pinned-original-source.tex; proof span is lines 2804-2856, the complete author proof (51 body lines).\par\medskip
\noindent\textit{Editorial note.} Navigation stubs: 11 cited-but-not-included labels are supplied as non-load-bearing pointers with their pinned tags where the snapshot tag index has them.\par\medskip
\noindent\textit{Editorial note.} extra\_packages is empty on purpose: the retained preamble already loads xy with options and a second declaration clashes in the assembler wrapper.\par\medskip
\noindent\textit{Editorial note.} Text is retained verbatim, including the author's own typos and placeholders.\par\medskip
\noindent\textit{Editorial note.} Queue index 3 of queue\_stacks-properties.json (pb=2019, stmt 2790-2802, proof 2804-2856).\par\medskip
\noindent\textit{Editorial note.} stacks-properties wave-1: fresh adjudication accept (2026-10-14); pool scanned while CONTROL.md is down (0 bytes) - §2 row deferred per Chief, claim reported to Chief 08:0x.\par\medskip
\noindent\textit{Allowed standard prerequisites (editorial).}\begin{itemize}
\item Explicit prerequisite boundary. The cited inputs are prerequisites; the argument in the retained span is not replaced by any cited result.
\end{itemize}\medskip
\ifdefined\hypertarget\hypertarget{sourcepointer-1}{}\fi
\noindent\textit{Source reference.} Original lemma 04Y5: sources/stacks-stacks-properties/pinned-original-source.tex, source lines 2856--2856. Cited by the retained statement or proof but not part of the collected spans; supplied only as a navigation pointer. The argument does not depend on its content.\par\medskip
\ifdefined\hypertarget\hypertarget{sourcepointer-2}{}\fi
\noindent\textit{Source reference.} Original lemma 03YQ: sources/stacks-stacks-properties/pinned-original-source.tex, source lines 2856--2856. Cited by the retained statement or proof but not part of the collected spans; supplied only as a navigation pointer. The argument does not depend on its content.\par\medskip
\ifdefined\hypertarget\hypertarget{sourcepointer-3}{}\fi
\noindent\textit{Source reference.} Original lemma 04T4: sources/stacks-stacks-properties/pinned-original-source.tex, source lines 2856--2856. Cited by the retained statement or proof but not part of the collected spans; supplied only as a navigation pointer. The argument does not depend on its content.\par\medskip
\ifdefined\hypertarget\hypertarget{sourcepointer-4}{}\fi
\noindent\textit{Source reference.} Original lemma 06FH: sources/stacks-stacks-properties/pinned-original-source.tex, source lines 2856--2856. Cited by the retained statement or proof but not part of the collected spans; supplied only as a navigation pointer. The argument does not depend on its content.\par\medskip
\ifdefined\hypertarget\hypertarget{sourcepointer-5}{}\fi
\noindent\textit{Source reference.} Original lemma 06FI: sources/stacks-stacks-properties/pinned-original-source.tex, source lines 2856--2856. Cited by the retained statement or proof but not part of the collected spans; supplied only as a navigation pointer. The argument does not depend on its content.\par\medskip
\ifdefined\hypertarget\hypertarget{sourcepointer-6}{}\fi
\noindent\textit{Source reference.} Original lemma lemma-monomorphism: sources/stacks-stacks-properties/pinned-original-source.tex, source lines 1674--1687. Cited by the retained statement or proof but not part of the collected spans; supplied only as a navigation pointer. The argument does not depend on its content.\par\medskip
\ifdefined\hypertarget\hypertarget{sourcepointer-7}{}\fi
\noindent\textit{Source reference.} Original lemma lemma-monomorphism-into-point: sources/stacks-stacks-properties/pinned-original-source.tex, source lines 2759--2764. Cited by the retained statement or proof but not part of the collected spans; supplied only as a navigation pointer. The argument does not depend on its content.\par\medskip
\ifdefined\hypertarget\hypertarget{sourcepointer-8}{}\fi
\noindent\textit{Source reference.} Original lemma lemma-points-cartesian: sources/stacks-stacks-properties/pinned-original-source.tex, source lines 1019--1036. Cited by the retained statement or proof but not part of the collected spans; supplied only as a navigation pointer. The argument does not depend on its content.\par\medskip
\ifdefined\hypertarget\hypertarget{sourcepointer-9}{}\fi
\noindent\textit{Source reference.} Original lemma lemma-unique-point: sources/stacks-stacks-properties/pinned-original-source.tex, source lines 2654--2662. Cited by the retained statement or proof but not part of the collected spans; supplied only as a navigation pointer. The argument does not depend on its content.\par\medskip
\ifdefined\hypertarget\hypertarget{sourcepointer-10}{}\fi
\noindent\textit{Source reference.} Original lemma lemma-unique-point-better: sources/stacks-stacks-properties/pinned-original-source.tex, source lines 2710--2717. Cited by the retained statement or proof but not part of the collected spans; supplied only as a navigation pointer. The argument does not depend on its content.\par\medskip
\ifdefined\hypertarget\hypertarget{sourcepointer-11}{}\fi
\noindent\textit{Source reference.} Original lemma 06G4: sources/stacks-stacks-properties/pinned-original-source.tex, source lines 2856--2856. Cited by the retained statement or proof but not part of the collected spans; supplied only as a navigation pointer. The argument does not depend on its content.\par\medskip
\begin{lemma}
\label{lemma-improve-unique-point}
Let $\mathcal{Z}$ be an algebraic stack. Assume $\mathcal{Z}$ satisfies
the equivalent conditions of
Lemma \ref{lemma-unique-point}.
Then there exists a unique strictly full subcategory
$\mathcal{Z}' \subset \mathcal{Z}$ such that
$\mathcal{Z}'$ is an algebraic stack which satisfies the equivalent
conditions of
Lemma \ref{lemma-unique-point-better}.
The inclusion morphism $\mathcal{Z}' \to \mathcal{Z}$ is a monomorphism
of algebraic stacks.
\end{lemma}

\begin{proof}
The last part is immediate from the first part and
Lemma \ref{lemma-monomorphism}.
Pick a field $k$ and a morphism $\Spec(k) \to \mathcal{Z}$
which is surjective, flat, and locally of finite type.
Set $U = \Spec(k)$ and $R = U \times_\mathcal{Z} U$.
The projections $s, t : R \to U$ are locally of finite type.
Since $U$ is the spectrum of a field, it follows that
$s, t$ are flat and locally of finite presentation (by
Morphisms of Spaces, Lemma
\ref{spaces-morphisms-lemma-noetherian-finite-type-finite-presentation}).
We see that $\mathcal{Z}' = [U/R]$ is an algebraic stack by
Criteria for Representability,
Theorem \ref{criteria-theorem-flat-groupoid-gives-algebraic-stack}.
By
Algebraic Stacks, Lemma \ref{algebraic-lemma-map-space-into-stack}
we obtain a canonical morphism
$$
f : \mathcal{Z}' \longrightarrow \mathcal{Z}
$$
which is fully faithful. Hence this morphism is representable by
algebraic spaces, see
Algebraic Stacks, Lemma
\ref{algebraic-lemma-characterize-representable-by-algebraic-spaces}
and a monomorphism, see
Lemma \ref{lemma-monomorphism}.
By
Criteria for Representability,
Lemma \ref{criteria-lemma-flat-quotient-flat-presentation}
the morphism $U \to \mathcal{Z}'$ is surjective, flat, and locally of finite
presentation. Hence $\mathcal{Z}'$ is an algebraic stack which satisfies
the equivalent conditions of
Lemma \ref{lemma-unique-point-better}.
By
Algebraic Stacks, Lemma \ref{algebraic-lemma-equivalent}
we may replace $\mathcal{Z}'$ by its essential image in $\mathcal{Z}$.
Hence we have proved all the assertions of the lemma except for the
uniqueness of $\mathcal{Z}' \subset \mathcal{Z}$. Suppose that
$\mathcal{Z}'' \subset \mathcal{Z}$ is a second such algebraic stack.
Then the projections
$$
\mathcal{Z}'
\longleftarrow
\mathcal{Z}' \times_\mathcal{Z} \mathcal{Z}''
\longrightarrow
\mathcal{Z}''
$$
are monomorphisms. The algebraic stack in the middle is nonempty by
Lemma \ref{lemma-points-cartesian}.
Hence the two projections are isomorphisms by
Lemma \ref{lemma-monomorphism-into-point}
and we win.
\end{proof}

\end{document}
